% Plantilla de actividad IMTC
\documentclass[spanish,letterpaper,oneside]{article}

\def\documenttitle {Introduccion A La Ciencia De Datos}
\def\documentsubtitle {Actividad 1 - introduccion-a-la-ciencia-de-datos-imtc}
\def\documentsubject {Ingenieria en Mecatronica}
\def\documentauthor {Martin Jonathan de la Cruz}
\def\coursename {Introduccion A La Ciencia De Datos}
\def\coursecode {IMTC}
\def\universityname {Universidad Autonoma de Nuevo Leon}
\def\universityfaculty {Facultad de Ingenieria Mecanica y Electrica}
\def\universitydepartment {\coursename}
\def\universitydepartmentimage {UANL/assets-uanl/logo-uanl.png}
\def\universitydepartmentimagecfg {width=3.2cm}
\def\universitylocation {San Nicolas de los Garza, Nuevo Leon}
\def\authortable {
  \begin{tabular}{ll}
    \textbf{Alumno:} & Martin Jonathan de la Cruz \\
    Matricula: & UANL-IMTC \\
    Programa: & Ingenieria en Mecatronica \\
    Asignatura: & \coursename \\
    Codigo: & \coursecode \\
    Figura docente: & Nombre por definir \\
    & \\
    \multicolumn{2}{l}{Fecha de realizacion: \today} \\
    \multicolumn{2}{l}{\universitylocation}
  \end{tabular}
}

\input{template}

\begin{document}
\templatePortrait
\templatePagecfg
\onehalfspacing

\section{Actividad 1}
\textbf{Consigna tecnica:} Documentar y resolver el reto academico de la semana.

\subsection{Desarrollo}
Incluir evidencias, analisis y referencias.

\subsection{Cierre}
Registrar resultados, limitaciones y mejoras.

% Ciclo PDF: bibliographystyle omitido; la plantilla AulaTeX ya define el estilo bibliográfico
\bibliography{introduccion-a-la-ciencia-de-datos-imtc}

% --- Ciclo A: refuerzo editorial materializado ---
\section*{Refuerzo editorial Ciclo A}
Este apartado consolida la mejora editorial incremental del nodo a partir del estándar maduro de AulaTeX. El refuerzo atiende brechas detectadas de bibliografía, citas, análisis propio, transferencia y cierre argumentado, sin sustituir la consigna original ni copiar contenido de los casos maestros.

\subsection*{Tesis de trabajo}
La actividad debe sostener una postura verificable: el producto solicitado no sólo organiza información, sino que permite interpretar el problema disciplinar, justificar decisiones y reconocer sus consecuencias académicas o profesionales.

\subsection*{Análisis propio}
El hallazgo central es que una entrega madura requiere pasar de la descripción a la interpretación. Por ello, el desarrollo debe explicar qué revela el producto construido, por qué es relevante para la materia y qué criterio permite evaluar su pertinencia. Esta operación convierte la evidencia reunida en argumento académico.

\subsection*{Transferencia profesional}
La transferencia consiste en vincular el producto con situaciones reales de desempeño: diagnóstico, toma de decisiones, comunicación de resultados, cumplimiento normativo, intervención educativa o resolución de problemas según el campo disciplinar. Así, la actividad adquiere valor formativo más allá de la plantilla.

\subsection*{Cierre argumentado}
En consecuencia, la calidad de la actividad depende de que el producto visible, las fuentes y la conclusión formen una secuencia coherente. La conclusión debe recuperar la tesis, indicar el principal aprendizaje y señalar una implicación práctica o conceptual.

% Brechas locales atendidas por Ciclo A: tesis, analisis_propio, transferencia, conclusion_argumentada, bibliografia_insuficiente, citas_insuficientes
% Reglas globales aplicadas:
% - Priorizar cierre de brecha recurrente: bibliografia_insuficiente (427 nodos).
% - Priorizar cierre de brecha recurrente: conclusion_argumentada (375 nodos).
% - Priorizar cierre de brecha recurrente: citas_insuficientes (328 nodos).
% - Priorizar cierre de brecha recurrente: analisis_propio (288 nodos).
% --- Fin Ciclo A ---


\subsection*{Citas de refuerzo Ciclo A}
La mejora editorial se vincula con la memoria documental disponible mediante las referencias \citep{programa-introduccion-a-la-ciencia-de-datos-imtc,cicloARefuerzo02,cicloARefuerzo03,cicloARefuerzo04,cicloARefuerzo05,cicloARefuerzo06}. Estas citas deben revisarse en una etapa disciplinar fina para sustituir o complementar fuentes genéricas por literatura específica del tema.

\end{document}
